
% ============================================================================
\clearpage
\section*{v4.10.0-r21 Source-Only: C5 Restated per Frame: Port Shares, the Frame-Free Floor and the Bridge Question}
\addcontentsline{toc}{section}{v4.10.0-r21 Source-Only: C5 Restated per Frame: Port Shares, the Frame-Free Floor and the Bridge Question}
\label{sec:v4100-r21-c5-restated}
% ============================================================================

\paragraph{Purpose.}
Layer v4.10.0-r20.4 found that C5 in its old form asks, in the measure of one frame, about a passage between two frames, and left the
restatement to the author.  The author has decided for it (2026-11-16).  This layer restates C5 in the way v4.9.0-r61 treated OP4: it
types the quantities by frame, splits the old entry \(C5_{\rm zero}\) of v4.7.0-r86 into three questions, and records the status of
each.  It also corrects one sentence of r20.4 and relates the collapse state of the main body to the two classes of fields.  Exploration, not blind: every computation is fixed analytically beforehand.
\textbf{The forbidden assertion \(C5=0\) is not made, and C5 stays open.}  The notes are
\texttt{notes/\allowbreak research/\allowbreak v4\_10\_0\_r21/} and
\texttt{notes/\allowbreak research/\allowbreak c5\_restatement\_2026\_11\_16/}.  Earlier layers are not rewritten.

\noindent\textbf{Source-only status.}
This v4.10.0-r21 layer is source-only.  It does not build a PDF, does not assign a new DOI,
and does not authorize a public release.  The v5.1.0 Version DOI
\href{https://doi.org/10.5281/zenodo.24237112}{10.5281/zenodo.24237112} and the Concept DOI
\href{https://doi.org/10.5281/zenodo.19210728}{10.5281/zenodo.19210728} are unchanged by it.

\subsection*{v4.10.0-r21 the restatement}
\label{subsec:v4100-r21-restatement}

\paragraph{Definition v4.10.0-r21.1 (shares and floor).}
On the doublet \((E,B)\) of v4.9.0-r04, with \(2u=E^2+B^2\), \(s=(E+iB)^2\) and \(b_x=2E\cdot B/2u\), put
\[
C5_M=\tfrac12(1+b_x),\qquad C5_T=\tfrac12(1-b_x),\qquad \phi=\frac{|s|}{2u}=|b| .
\]
\(C5_M\) is the share of the content in the compensator port \(C_{\rm HC}=\tfrac12(1+\sigma_x)\), read by the matter frame; \(C5_T\) is the
share in the partner port, which the tachyon frame reads (v4.9.0-r28); \(\phi\) is the floor of v4.9.0-r10.2(b) relative to the content.
For the selective collar recursion the cumulative shares are \(C5_M^{(n)}=(3/4)^n\) and \(C5_T^{(n)}=1-(3/4)^n\) (r20.4(b)).

\paragraph{Restatement v4.10.0-r21.2 (three questions in place of \(C5_{\rm zero}\)).}
\begin{enumerate}[leftmargin=2em]
\item[\(C5_{\rm tr}\)] \emph{Transfer question, typed by frame.}  Is there a finite stage with \(C5_M=0\), equivalently \(C5_T=1\)?
\item[\(C5_P\)] \emph{Photon-point question, frame-free.}  Does a non-null state with non-zero content reach \(\phi=0\) at a finite value of the
parameter?
\item[\(C5_{\rm br}\)] \emph{Bridge question.}  Is reaching \(\phi=0\) a discharge of the obstruction in the sense of v4.7.0-r86 and r88?
\end{enumerate}
The old entry ``\(\nexists N<\infty:\Phi_5^{(N)}=0\)'' read the score of one frame and is \(C5_{\rm tr}\); the route of r88, shape (iii) with
requirement (\(\alpha\)), concerns the obstruction and is \(C5_P\) together with \(C5_{\rm br}\).

\subsection*{v4.10.0-r21 results}
\label{subsec:v4100-r21-results}

\paragraph{Theoremlet v4.10.0-r21.3 (shares, exchange, floor).}
\begin{enumerate}[leftmargin=2em]
\item[(a)] \(C5_M+C5_T=1\) for every field, and for the cumulative shares at every stage.
\item[(b)] The tachyon map \(z\mapsto iz\) exchanges \(C5_M\) and \(C5_T\); the floor \(\phi\) is fixed by \(R\), \(T\) and \(J_B\).
\item[(c)] \(C5_M=C5_T=\tfrac12\) holds exactly for \(E\cdot B=0\).  This is the whole matter--tachyon axis of the disk, including a pure
electric field (\(\phi=1\)).  The photon point \(P\) is \(\phi=0\); equal shares are necessary for it and not sufficient.
\end{enumerate}

\paragraph{Correction to v4.10.0-r20.4.}
The reading of r20.4 named \(P\) as the frame-free point of a statement with the shares \(C5_M\), \(C5_T\), and the wiki called it the point
where both shares are \(\tfrac12\).  By r21.3(c) equal shares do not characterise \(P\).  The frame-free quantity is the floor \(\phi\), and
\(P\) is its zero.  The results r20.1 to r20.4(a)--(c) are unaffected.

\paragraph{Theoremlet v4.10.0-r21.4 (the transfer question is a relocation; conditional on the reading of v4.9.0-r10).}
\begin{enumerate}[leftmargin=2em]
\item[(a)] \(C5_M=0\) holds exactly for \(E=-B\).  Such fields exist statically; they have \(C5_T=1\) and the maximal floor \(\phi=1\).
\item[(b)] By r21.3(a) an empty compensator port is a full partner port.  In the typed question nothing vanishes; the content changes
frame.  This is v4.9.0-r10.2(a) in typed form, and the passage falls under requirement (\(\delta\)) of r88.
\item[(c)] Inside the selective collar recursion \(C5_M^{(n)}>0\) for every finite \(n\) (v4.7.0-r86.2).
\end{enumerate}
Reading: for the transports of Reading~B, \(C5_{\rm tr}\) asks for a complete change of frame and not for a vanishing.  It carries no
discharge.

\paragraph{Theoremlet v4.10.0-r21.5 (the photon-point question).}
\begin{enumerate}[leftmargin=2em]
\item[(a)] Boosts, spatial rotations and duality rotations keep \(|s|\) (v4.9.0-r10.2, r19.1); the non-selective flow of v4.9.0-r27 reaches
\(P\) only in the limit, except at \(\alpha=\pi/4\).
\item[(b)] Free time evolution reaches \(\phi=0\) at a finite time for the codimension-one family of r20.1.
\item[(c)] The cumulative shares of the collar recursion are equal at no integer stage: \((3/4)^n=\tfrac12\) at \(n=2.409\), with
\((3/4)^2=0.5625\) and \((3/4)^3=0.4219\).
\end{enumerate}
\(C5_P\) thus has a positive answer as a statement about fields.  It is not a statement about C5.

\paragraph{Theoremlet v4.10.0-r21.6 (a lasting zero).}
For a time-harmonic field the mean of \(|s(t)|^2\) over a period is \(\tfrac1{16}\bigl(|s_+|^2+4|w|^2+|s_-|^2\bigr)\).  Hence \(s(t)=0\) for all \(t\)
if and only if \(s_+=s_-=w=0\).  Free time evolution keeps \(|s_+|\), \(|s_-|\) and \(|w|\).  A field that stays at the photon point was there
from the start (plane waves, circular and linear); free time evolution leads into no lasting zero.  The example of r20.1 has mean
\(0.305\).

\paragraph{Proof sketch.}
r21.3 and r21.4(a) from the definitions; \(b_x=-1\) forces \(2E\cdot B=-(E^2+B^2)\), that is \(|E+B|^2=0\).  r21.6 by Parseval on
\(s(t)=\tfrac14\bigl(s_+e^{-2i\omega t}+2w+\bar s_-e^{2i\omega t}\bigr)\) (r16.8).  All items were checked in
\texttt{restatement\_\allowbreak checks.py}: \(1000\) fields for r21.3, \(300\) fields for r21.6 (relative deviation \(9\cdot10^{-16}\)).  \(\square\)

\paragraph{Remark v4.10.0-r21.7 (where the open question now sits).}
After the restatement the open part of C5 is \(C5_{\rm br}\) alone.  It depends on three points.
\begin{enumerate}[leftmargin=2em]
\item[(B1)] The reading of the obstruction as the floor of the instantaneous field (v4.9.0-r10, r20.2).  It remains a reading.
\item[(B2)] Instant or duration.  v4.7.0-r86 asks for a finite stage \(N\) with \(\Phi_5^{(N)}=0\) and does not say whether the zero must
persist.  Under an \emph{event reading} an invariant instant suffices, and the null instant of r20.1 is a candidate, subject to (B1)
and (B3).  Under a \emph{duration reading} the zero must last; then by r21.6 free time evolution supplies no witness and the
candidate of r20 drops out.
\item[(B3)] Presence.  Why fields of the codimension-one family should be present is open (r20).
\end{enumerate}
(B2) is the author's decision.  On 2026-11-16 the author set the event reading as a working assumption (r21.10); it is set, not derived.

\paragraph{Theoremlet v4.10.0-r21.8 (two classes).}
The exploration k220 of 2026-11-07 found the orbits of \(SO(3,\mathbb C)\times U(1)_{\rm duality}\) on \(z=E+iB\): the shells \(|s|=r>0\), the
punctured null cone and \(\{0\}\).  Add the scale, \(z\mapsto\lambda z\) with \(\lambda\in\mathbb C\setminus\{0\}\) (duality phase and dilation).
\begin{enumerate}[leftmargin=2em]
\item[(a)] All non-null fields form one orbit: for \(s_1,s_2\ne0\) there are \(\lambda\) and \(L\in SO(3,\mathbb C)\) with \(z_2=\lambda Lz_1\).
\item[(b)] The non-zero null fields form one orbit of \(SO(3,\mathbb C)\) alone.
\item[(c)] Each class has a distinguished representative.  A non-null field is taken by a finite boost to the frame \(E\parallel B\), where
\(\phi=1\); this is the pure state, the reference frame R of v4.9.0-r24.3.  A null field has \(\phi=0\) in every frame; this is the photon
point \(P\).  Under boosts \(\phi\) takes every value in \((0,1]\) on the first class and only \(0\) on the second.
\item[(d)] The rapidity of the boost to R obeys \(\phi=\operatorname{sech}2\eta_R\).  Along the null instant of r20.1 it diverges
(\(\eta_R=0.74,\ 1.33,\ 2.02,\ 2.83,\ 3.98\) at \(t=-0.6,\ -0.2,\ -0.05,\ -0.01,\ -0.001\)).  Free time evolution reaches at a finite time what
boosts reach only at infinite rapidity.
\end{enumerate}
Reading: below every frame readout there is one predicate left, whether the field lies in the null class or not.  The shares, the
score \((3/4)^n\) and the value of the floor are readouts of a frame above it.  Checked in \texttt{two\_\allowbreak classes.py}
(\(2000\) pairs each for (a) and (b), \(1000\) fields for (c)).

\paragraph{Remark v4.10.0-r21.9 (the collapse state of the main body in present terms; reading).}
The main body calls \(\zminus=0\) the collapse state, the zero of the frame variable, and identifies it with the photon and with the
unique frame-invariant object (Theorem~\ref{thm:unity}, Theorem~\ref{thm:origin}).  v4.9.0-r24.3 typed the zero of the offset as R, the
pure state, and separated it from the photon \(P\) of Convention~F.  With r21.8 the two roles of the collapse state are the two
distinguished representatives:
\begin{center}
\small
\begin{tabular}{@{}>{\raggedright\arraybackslash}p{0.40\textwidth}>{\raggedright\arraybackslash}p{0.52\textwidth}@{}}
\toprule
role of \(\zminus=0\) in the main body & present description\\
\midrule
zero of the offset; target of the \(\gL^{-k}\) conversions; \(\Lambda=0\) & R: \(\phi=1\), \(\eta_R=0\); a frame, reached by a finite boost (r24.3)\\
the invariant object; the photon & \(P\): \(\phi=0\), the null class, fixed point of duality (k220); \(\eta_R=\infty\)\\
``we can never fully reach \(\zminus=0\)'' (Theorem~\ref{thm:dark-offset}, v4.8.0-r18) & as a statement of reachability it fits \(P\) (no transport, boosts only in the limit), not R\\
compensator output as projection onto the collapse state (Theorem~\ref{thm:unity}(iv)) & neither: the image is not null (v4.10.0-r11.2(f)); a projection loses information (k220)\\
\bottomrule
\end{tabular}
\end{center}
R and \(P\) are the two ends of one rapidity line, \(\eta_R=0\) and \(\eta_R=\infty\); the main body, which works in one frame (r24), wrote both
as \(\zminus=0\).  The third row is a tension with the typing of r24.3 and is recorded, not resolved.  The main body is not rewritten.

\paragraph{Remark v4.10.0-r21.10 (working assumption for (B2)).}
If the lowest level carries only the class predicate of r21.8, membership at an instant is the primitive notion: it is
Lorentz-invariant (r20.1(c)), while a duration needs a time parameter that the class level does not have.  On this ground the author
set the event reading of (B2) as a working assumption on 2026-11-16.  It is a setting.  That the scale and the duality phase are
divided out at the lowest level is a setting as well; both are symmetries of the free field only (k220).  (B1) and (B3) stay open.

\paragraph{What this does not give.}
No C5 discharge.  The classification of \(C5_{\rm tr}\) holds within the reading of v4.9.0-r10 and for the transports of Reading~B.
\(C5_P\) is answered for fields only.  \(C5_{\rm br}\) is open.

\subsection*{v4.10.0-r21 status in the three frames}
\label{subsec:v4100-r21-frames}

\begin{center}
\small
\begin{tabular}{@{}>{\raggedright\arraybackslash}p{0.30\textwidth}>{\raggedright\arraybackslash}p{0.62\textwidth}@{}}
\toprule
quantity or question & M / T / P\\
\midrule
\(C5_M\), \(C5_T\) & M reads \(C5_M\), T reads \(C5_T\); the tachyon map exchanges them\\
floor \(\phi\) & frame-free; \(P\) is \(\phi=0\)\\
\(C5_{\rm tr}\) & a passage from M to T\\
\(C5_P\), \(C5_{\rm br}\) & frame-free\\
class predicate (r21.8) & frame-free; R is a frame of the non-null class, \(P\) is the null class\\
\bottomrule
\end{tabular}
\end{center}

\subsection*{v4.10.0-r21 register entries}
\label{subsec:v4100-r21-register}

\begin{center}
\small
\begin{tabular}{@{}>{\raggedright\arraybackslash}p{0.34\textwidth}>{\raggedright\arraybackslash}p{0.58\textwidth}@{}}
\toprule
statement & status after v4.10.0-r21\\
\midrule
C5 in the old form & mixes frames (r20.4); restated as \(C5_{\rm tr}\), \(C5_P\), \(C5_{\rm br}\) (r21.2)\\
\(C5_{\rm tr}\) (transfer, typed) & a relocation, not a vanishing, under the r10 reading; excluded inside the recursion (r21.4)\\
\(C5_P\) (photon point, frame-free) & reached at finite time by free evolution on a codimension-one family; a statement about fields (r21.5)\\
lasting zero of \(s(t)\) & only for \(s_+=s_-=w=0\), present from the start (r21.6)\\
\(C5_{\rm br}\) (bridge) & open; depends on (B1), (B2), (B3); (B2) is the author's decision (r21.7)\\
equal shares and \(P\) & necessary, not sufficient; correction to r20.4 (r21.3(c))\\
classes of non-zero fields & two under Lorentz group, duality and scale: null or not; representatives \(P\) and R (r21.8)\\
collapse state \(\zminus=0\) of the main body & two roles: zero of the offset (R) and invariant object (\(P\)); reading, one tension recorded (r21.9)\\
(B2) instant or duration & event reading set as working assumption by the author; not derived (r21.10)\\
\(C5_{\rm zero}\) & open; no discharge, no assertion \(C5=0\)\\
\bottomrule
\end{tabular}
\end{center}

\subsection*{v4.10.0-r21 ledger and guard update}
\label{subsec:v4100-r21-ledger-guard}

\[
R_{\rm chk}^{v4100\text{-}r21}=(1,1,0)
\]
for (restatement recorded; computations checked in the cited script; the bridge question open).

The layer spends no guard:
\[
R_{\rm glob}^{(v4100\text{-}r21,+)}=(0,1,1,1,1),
\qquad
\Gamma_{\rm guard}^{v4100\text{-}r21}=(1,1,1,1,1),
\qquad
\beta_4=1/256\ \text{unspent}.
\]
The layer does not imply exact finite \(C5\) zero closure, exact all-readout terminality, \(C6/T2\), a global photon-clock/RH theorem,
or global all-readout closure.
